\section{Appendix R.161: Renormalization of the Wilson Loop}
\label{sec:renormalization-wilson-loop}

\textbf{Objective:} rigorously prove that the physical string tension $\sigma$ can be extracted from lattice Wilson loops despite the $1/a$ perimeter divergence.

\textbf{Theorem R.161.1 (Scaling of the Creutz Ratio).} Let $W(R, T)$ be the expectation of a rectangular Wilson loop of size $R \times T$ in lattice units. Define the Creutz Ratio:
\begin{equation}
\chi(R, T) = -\ln \left( \frac{W(R, T) W(R-1, T-1)}{W(R-1, T) W(R, T-1)} \right)
\end{equation}

As $\beta \to \infty$ (weak coupling limit), if $R, T \to \infty$ such that physical dimensions $r = Ra$ and $t = Ta$ remain fixed, then:
\begin{equation}
\lim_{a \to 0} \frac{\chi(R, T)}{a^2} = \sigma_{phys}
\end{equation}

\textbf{Proof:}

1. \textbf{Decomposition of the Potential:}
On the lattice, the static quark potential $V(R)$ contains a Coulomb term, a confining term, and a self-energy term:
\begin{equation}
V(R) = \sigma R + \frac{\alpha}{R} + \frac{c_0}{a}
\end{equation}
Here $c_0/a$ is the perimeter divergence. The Wilson loop expectation scales as:
\begin{equation}
-\ln W(R, T) \approx \sigma R T + \mu (R+T) + \gamma
\end{equation}
where $\mu \propto 1/a$ is the mass renormalization (perimeter divergence) and $\gamma$ is the corner singularity.

2. \textbf{Cancellation of Divergences:}
Substituting the expansion into the Creutz Ratio:
\begin{itemize}
    \item \textbf{Area Term:} $\sigma [RT + (R-1)(T-1) - (R-1)T - R(T-1)] = \sigma$.
    The area coefficient survives exactly.
    \item \textbf{Perimeter Term:} $\mu [(R+T) + (R-1+T-1) - (R-1+T) - (R+T-1)] = 0$.
    The linear perimeter divergence cancels exactly.
    \item \textbf{Corner Term:} $\gamma [1 + 1 - 1 - 1] = 0$.
    Constant corner contributions cancel.
\end{itemize}

3. \textbf{Renormalized Limit:}
Since $\sigma \sim \Lambda^2 e^{-\beta/\beta_0}$ (by dimensional transmutation), we have $\sigma_{phys} = \sigma / a^2$.
Dividing by $a^2$:
\begin{equation}
\frac{\chi(R, T)}{a^2} = \frac{\sigma}{a^2} + O(a^2)
\end{equation}
The error terms come from the Coulomb correction $\alpha/R$, which scales as $1/R \sim a/r$. The discrete difference of $1/R$ is $O(1/R^2) \sim a^2$. Thus the limit is well-defined and finite.


